%%%%%%%%%%%%%%%%%%%%%%%%%%%%%%%%%%%%%%%%%
% "ModernCV" CV and Cover Letter
% LaTeX Template
% Version 1.11 (19/6/14)
%
% This template has been downloaded from:
% http://www.LaTeXTemplates.com
%
% Original author:
% Xavier Danaux (xdanaux@gmail.com)
%
% License:
% CC BY-NC-SA 3.0 (http://  creativecommons.org/licenses/by-nc-sa/3.0/)
%
% Important note:
% This template requires the moderncv.cls and .sty files to be in the same 
% directory as this .tex file. These files provide the resume style and themes 
% used for structuring the document.
%
%%%%%%%%%%%%%%%%%%%%%%%%%%%%%%%%%%%%%%%%%

%----------------------------------------------------------------------------------------
%	PACKAGES AND OTHER DOCUMENT CONFIGURATIONS
%----------------------------------------------------------------------------------------

\documentclass[11pt,a4paper,sans]{moderncv} % Font sizes: 10, 11, or 12; paper sizes: a4paper, letterpaper, a5paper, legalpaper, executivepaper or landscape; font families: sans or roman

\moderncvstyle{banking} % CV theme - options include: 'casual' (default), 'classic', 'oldstyle' and 'banking'
\moderncvcolor{blue} % CV color - options include: 'blue' (default), 'orange', 'green', 'red', 'purple', 'grey' and 'black'

  
\usepackage[T1]{fontenc}
\usepackage[utf8]{inputenc}
\usepackage{lmodern}
\usepackage{needspace}
\usepackage[scale=0.90]{geometry} % Reduce document margins
\setlength{\hintscolumnwidth}{2.3cm} %
% personal data
\name{Michael}{Arbel}
\title{Research Scientist in Machine Learning}
\address{Inria Grenoble -- Rhône-Alpes}{Grenoble}{France}
\email{michael.n.arbel@gmail.com}                               % optional, remove / comment the line if not wanted
\homepage{michaelarbel.github.io}

% Social icons
\social[linkedin]{michael-arbel-0a38a655}
\social[github]{MichaelArbel}
\usepackage[defernumbers=true,style=numeric,maxbibnames=9,maxcitenames=1,uniquelist=true,backend=biber,giveninits=true,sorting=ydnt,doi=true,url=false]{biblatex} 
%\addbibresource{https://raw.githubusercontent.com/MichaelArbel/MichaelArbel.github.io/refs/heads/master/assets/bibliography/bibliography.bib}
\addbibresource{../assets/bibliography/bibliography.bib}

%%%%%%%%%%%%%%%%%%%%%%%%%%%%%%%%%%%%%%%%%%%%%%%%%%%%%%%%%%%%%%%%%%%%%%%%%%%%%%%%%%%%%%%%%%%%%%%%%%%%%%%%%%%%%%%%%%%%%%%%%%%%%%%%%%%%%%%%%%%%%%%%%%%%%%%%%%%%
%%%%%%%%%%%%%%%%%%%%%%%%%%%%%%%%%%%%%%%%%%%%%%%%%%%%%%%%%%%%%%%%%%%%%%%%%%%%%%%%%%%%%%%%%%%%%%%%%%%%%%%%%%%%%%%%%%%%%%%%%%%%%%%%%%%%%%%%%%%%%%%%%%%%%%%%%%%%
% END BLACK MAGIC
%%%%%%%%%%%%%%%%%%%%%%%%%%%%%%%%%%%%%%%%%%%%%%%%%%%%%%%%%%%%%%%%%%%%%%%%%%%%%%%%%%%%%%%%%%%%%%%%%%%%%%%%%%%%%%%%%%%%%%%%%%%%%%%%%%%%%%%%%%%%%%%%%%%%%%%%%%%%
%%%%%%%%%%%%%%%%%%%%%%%%%%%%%%%%%%%%%%%%%%%%%%%%%%%%%%%%%%%%%%%%%%%%%%%%%%%%%%%%%%%%%%%%%%%%%%%%%%%%%%%%%%%%%%%%%%%%%%%%%%%%%%%%%%%%%%%%%%%%%%%%%%%%%%%%%%%%

% 


\begin{document}

\makecvtitle % Print the CV title

\section{Research profile}
\cvitem{}{I develop principled machine learning methods that combine data with structure, simulators, and implicit constraints for reliable scientific inference and decision-making. My research connects statistical learning theory, optimization, and empirical work, with current directions in functional bilevel optimization, predictive representation learning, and reliable tabular learning.}

%----------------------------------------------------------------------------------------
%	EDUCATION SECTION
%----------------------------------------------------------------------------------------

\section{Academic positions}
\cventry{Since 2022}{THOTH team}{Research Scientist (Chargé de recherche), Inria}{Grenoble, France}{}{}
\cventry{2021-2022}{Working with Julien Mairal on bilevel optimization}{Starting Research Fellow at Inria}{Grenoble, France}{}{}

\section{Education}
\cventry{2016-2021}{Under the supervision of Arthur Gretton}{PhD in Machine Learning, University College London}{London, UK}{\textit{Pass with no revisions}}{
Dissertation: Regularization and Optimization of Generative Adversarial Networks.} 
\cventry{2014-2015}{Machine learning and computer vision (MVA)}{M.S. in applied mathematics at École Normale Supérieure}{Cachan, France}{\textit{High honors}}{Master's thesis at École Normale Supérieure de Paris under the supervision of  Stéphane Mallat, 
\\Subject: Wavelet analysis of long-range dependencies in simple grammars.}
\cventry{2011-2014}{Minor in physics}{B.S. \& M.S. in applied mathematics at École Polytechnique}{Palaiseau, France}{ GPA: 3.93/4}{
Master thesis at Princeton University, under the supervision of René Carmona,\\
Subject: Mean field games with a dependence on the distribution of the  control.}

\section{Industrial positions}
\cventry{2015-2016}{Computer Vision for neuromorphic cameras}{Research Engineer at Prophesee}{Paris, France}{}{}

%----------------------------------------------------------------------------------------
%	WORK EXPERIENCE SECTION
%----------------------------------------------------------------------------------------

\section{Grants and scholarships}
\cvitem{2027-2032}{\textbf{DELPHI, ERC Starting Grant}, \textbf{1.5M euros}. Awarded in 2026; scheduled to begin in October 2027. Implicitly constrained learning for scientific discovery and decision-making.}
\cvitem{2024-2028}{\textbf{BONSAI, ANR JCJC}, \textbf{374K euros}. Project coordinator. Bilevel Optimization for Simulation-based Inference; March 2024--February 2028.}
\cvitem{2021-2022}{Starting research grant from INRIA, \textbf{ 10K euros}. }
\cvitem{2016-2021}{Scholarship from Gatsby Computational Neuroscience Unit for the PhD program. 22K pounds/year}
\cvitem{2011-2014}{Scholarship from the Eiffel Excellence Scholarship program. 14K euros/year. }


\needspace{10\baselineskip}
\section{Honors and Awards}
\cvitem{2024}{Spotlight presentation at NeurIPS 2024, \emph{Functional Bilevel Optimization for Machine Learning}.}
\cvitem{2022}{Long oral presentation at AISTATS 2022.}
\cvitem{2021}{Long oral presentation at ICML 2021.}
\cvitem{2020}{Spotlight presentation at ICLR 2020.}
\cvitem{2018}{Best Poster Award at MSR AI Summer School, Cambridge.}
\cvitem{2014}{Award of the Financial Risk Chair of Polytechnique for the Master's thesis.}

\section{Services to the Community}
\subsection{Editorial Activities}
\cventry{2023}{Area chair}{ICLR 2023 and NeurIPS 2023}{}{}{}
\cventry{2020-present}{Board of reviewers}{Journal of Machine Learning Research (JMLR)}{}{}{}
\cventry{2018-2021}{Reviewer}{NeurIPS (2018-2021), ICLR (2019-2021), ICML (2021).}{}{}{}
\subsection{Organization of scientific events}
\cventry{2017-2019}{DeepMind/CSML Seminars}{UCL}{London, UK}{}{
Weekly research seminars in Machine Learning with invited speakers from the UK and Europe. Annual budget: 5K pounds.
}






\section{Research supervision}
\cvitem{Current PhDs}{Supervision or co-supervision of Kenta Vert (2024--ongoing), Eloïse Touron (2024--ongoing), Fares El Khoury (2024--ongoing), and Pierre-Louis Ruhlmann (2023--ongoing).}
\cvitem{PhD alumni}{Ieva Petrulionyte (2023--2026), co-supervised with Julien Mairal; Juliette Marrie (2022--2025), co-supervised with Julien Mairal and Diane Larlus.}
\cvitem{Master's}{Tristan Martin (2026); Eloïse Touron and Camille Touron (2024); Ieva Petrulionyte (2023).}
\cvitem{}{Research topics, collaborators, and representative student publications: \href{https://michaelarbel.github.io/team/}{michaelarbel.github.io/team/}.}

\section{Teaching graduate courses}
\cventry{2026-2027}{Co-instructor, \href{https://kernel-learning.github.io/}{From Basic Machine Learning Models to Advanced Kernel Learning}}{Université Grenoble Alpes}{Grenoble, France}{}{}
\cventry{Spring 2026}{Co-instructor, \href{https://mva-kernel-methods.github.io/course-page/}{Machine Learning with Kernel Methods}}{Joint IASD, MVA, and MASH course}{Dauphine--PSL / ENS Paris-Saclay}{}{}
\cvitem{Earlier teaching}{Kernel methods at ENS Paris-Saclay and advanced kernel learning at Université Grenoble Alpes since 2022; Intelligent Systems at Université Grenoble Alpes (2023--2024).}

\section{Selected research software}
\cvitem{\href{https://github.com/inria-thoth/mlxp}{MLXP}}{Main author. Python framework for reproducible machine learning experiments: launching, logging, and querying results. ACM REP 2024.}
\cvitem{\href{https://github.com/MichaelArbel/EquiTabPFN}{EquiTabPFN}}{Target-permutation equivariant prior-fitted network for tabular learning. NeurIPS 2025.}
\cvitem{\href{https://github.com/inria-thoth/funcBO}{funcBO}}{Functional bilevel optimization algorithms and applications. NeurIPS 2024 (Spotlight).}
\cvitem{\href{https://github.com/MichaelArbel/BGS-opt}{BGS-opt}}{Research code for AmIGO and Bilevel Games with Selection, including hyperparameter optimization and dataset distillation. ICLR 2022; NeurIPS 2022.}


\section{Selected Invited Talks}

\subsection{Upcoming invitations}
\cvitem{19 Oct. 2026}{\emph{When Predictive Self-Supervised Learning Becomes Canonical Correlation Analysis}. Gatsby Computational Neuroscience Unit, London.}
\cvitem{13 Oct. 2026}{\emph{What Structure Does Predictive Self-Supervised Learning Recover?} Institute for Mathematics of Toulouse.}

\subsection{Recent presentations}
\cvitem{29 June 2026}{\emph{Bilevel optimization: A functional perspective with application to self-supervised learning}. SciCADE, Edinburgh; minisymposium on bilevel optimization for scientific computing.}
\cvitem{17 June 2026}{\emph{PEIRA: Learning Predictive Encoders through Inter-View Regressor Alignment}. Optimization and Learning Day, LJK, Grenoble.}
\cvitem{1 June 2026}{\emph{Bilevel optimization: A functional perspective with application to self-supervised learning}. SIAM Conference on Optimization.}
\cvitem{21 Apr. 2026}{\emph{Learning Predictive Embeddings through Inter-View Regressor Alignment}. CILVR seminar, Center for Data Science, New York University.}
\cvitem{Mar. 2026}{\emph{A functional approach to differentiable programming in machine learning}. Workshop on Flows on Measure Spaces and Applications in Machine Learning, Oberwolfach.}
\cvitem{Feb. 2025}{\emph{Functional bilevel optimization: theory and algorithms}. RKHS Seminars, METU.}

\subsection{Invited talks at conferences and workshops}

\cvitem{2023}{Bayes Comp. Levi, Finland.}

\cvitem{2023}{Conference in Mathematics and Image Analysis. Berlin, Germany.}

\cvitem{2022}{Learning and Optimization in Luminy workshop. Luminy, France.}
\cvitem{2022}{ELISE Theory Workshop on Machine Learning Fundamentals. EURECOM, Sophia Antipolis, France.}
\cvitem{2020}{Workshop on Functional Inference and Machine Intelligence, EURECOM.}
\cvitem{2019}{Amazon Research Days (Berlin, Germany).}
\cvitem{2019}{Workshop on Recent developments in kernel methods, 2019, UCL (London, UK).}
\cvitem{2019}{Deep Learning Theory Kickoff Meeting 2019, MPI (Leipzig, Germany).}
\cvitem{2018}{Cambridge-Tübingen workshop 2018 (Tenerife, Spain).}


\subsection{Seminars}

\cvitem{2021}{NYU Center for Data Science. (NY, USA) Remote.}

\cvitem{2021}{Instituto Superior Técnico, Lisbon, Portugal (Remote).}

\cvitem{2020}{The Alan Turing Institute (London, UK).}
\cvitem{2020}{Department of Statistics, University of Oxford (Oxford, UK).}
\cvitem{2019}{The Alan Turing Institute (London, UK).}
\cvitem{2018}{Google Developer Group Reading and Thames Valley (Reading, UK).}

%----------------------------------------------------------------------------------------
%	COMPUTER SKILLS SECTION
%----------------------------------------------------------------------------------------

\section{Publications}
\nocite{*}
\printbibliography[heading=none,title=none]

%
%----------------------------------------------------------------------------------------
%	COMMUNICATION SKILLS SECTION
%----------------------------------------------------------------------------------------
%----------------------------------------------------------------------------------------
%	LANGUAGES SECTION
%----------------------------------------------------------------------------------------  

%----------------------------------------------------------------------------------------
%	INTERESTS SECTION
%----------------------------------------------------------------------------------------

%----------------------------------------------------------------------------------------
 
\end{document}
